%%%%%%%%%%  scan idx 567  |  folio 554  %%%%%%%%%%
\SGAPrintPage{554}%

\sgaStableTarget{sga6.e10.prop.6-2-2}\textbf{Proposition 6.2.2.} \uline{For every positive section} \(\underline{n}\) \uline{of} \(\mathbb Z\)\uline{, the étale sheaf of groups} \(\underline{\Der}(\Gl(\underline{n}))\) \uline{of commutators of} \(\Gl(\underline{n})\) \uline{is represented by}
\(\operatorname{Sl}(\underline{n})=\Ker(\det:\Gl(\underline{n})\to{G}_m)\).

\sgaStableTarget{sga6.e10.cor.6-2-3}\textbf{Corollary 6.2.3.} \uline{Every representation}
\sgaStableTarget{sga6.e10.eq.6-2-4}\begin{equation}\tag{6.2.4}
 \alpha:G=\Gl(\underline{n}_1)\times\cdots\times\Gl(\underline{n}_r)\longrightarrow{G}_m
\end{equation}
\uline{is of the form}
\sgaStableTarget{sga6.e10.eq.6-2-5}\begin{equation}\tag{6.2.5}
 \alpha(Y)(s)=(\det(s_1))^{\underline{m}_1}\cdots(\det(s_r))^{\underline{m}_r},
\end{equation}
\uline{where} \(Y\in\operatorname{Ob}(\Sch/S)\), \(s=(s_1,\ldots,s_r)\in G(Y)\) \uline{and} \(\underline{m}_1,\ldots,\underline{m}_r\) \uline{are sections of} \(\mathbb Z\) \uline{uniquely determined by} \(\alpha\).

\sgaStableTarget{sga6.e10.cor.6-2-6}\textbf{Corollary 6.2.6.} \uline{With the notation of} 6.1. \uline{one has}
\sgaStableTarget{sga6.e10.eq.6-2-6}\[
 \det\bigl(T_\alpha(\underline{E}_1,\ldots,\underline{E}_r)\bigr)
 =
 \det(\underline{E}_1)^{\otimes \underline{m}_1}\otimes\cdots\otimes
 \det(\underline{E}_r)^{\otimes \underline{m}_r},
\]
\uline{where} \(\underline{m}_1,\ldots,\underline{m}_r\) \uline{are sections of} \(\mathbb Z\) \uline{uniquely determined by} \(\alpha\).

\sgaStableTarget{sga6.e10.rem.6-2-7}\textbf{Remark 6.2.7.} One can determine the exponents \(\underline{m}_1,\ldots,\underline{m}_r\) in (6.2.5) and (6.2.6) by successively substituting in (6.2.5)
\(s_i=\lambda.\ \operatorname{id}.(\lambda\in{G}_m(Y))\), \(s_j=\operatorname{id}_{\Gl(\underline{n}_j)(Y)}\), for \(j\ne i\) \((i,j=1,\ldots,r)\), whence
\[
 \det(s_i)=\lambda^{\underline{n}_i}\quad\text{and}\quad
 \alpha(Y)(s)=\lambda^{\underline{n}_im_i}.
\]
For example, one thus deduces from (6.2.6)
\sgaStableTarget{sga6.e10.eq.6-2-8}\sgaStableTarget{sga6.e10.eq.6-2-9}\begin{align}
 \det(\underline{E}_1\otimes \underline{E}_2)&=(\det(\underline{E}_1))^{\otimes \underline{n}_2}\otimes(\det(\underline{E}_2))^{\otimes \underline{n}_1}, \tag{6.2.8}\\
 \det(\underline{E}_1\oplus \underline{E}_2)&=(\det(\underline{E}_1))\otimes(\det(\underline{E}_2)) \tag{6.2.9}
\end{align}
and
\sgaStableTarget{sga6.e10.eq.6-2-10}\begin{equation}\tag{6.2.10}
 \det\bigl(\bigwedge\limits^{\underline{m}} \underline{E}_1\bigr)=(\det(\underline{E}_1))^{\binom{\underline{n}_1-1}{\underline{m}-1}}
\end{equation}
etc.

%%%%%%%%%%  scan idx 568  |  folio 555  %%%%%%%%%%
\SGAPrintPage{555}%

\sgaStableTarget{sga6.e10.sec.6.3}\subsection*{6.3.}
Let us first show how 6.2.3 follows from 6.2.2. For this, it plainly suffices to prove (6.2.5) for the
\(s_{(i)}=(1,\ldots,s_i,\ldots,1)\) \((i=1,\ldots,r)\), so one may suppose \(r=1\), \(G=\operatorname{G\ell}(\underline{n}_1)\). By 6.2.2, \(\alpha\) factors through \(\operatorname{G\ell}(\underline{n}_1)/\operatorname{Sl}(\underline{n}_1)\), i.e., there is a commutative diagram
\[
\begin{tikzcd}[column sep=large,row sep=large]
\operatorname{G\ell}(\underline{n}_1) \arrow[rr,"\alpha"] \arrow[dr,"j"'] && {G}_m\\
& \operatorname{G\ell}(\underline{n}_1)/\operatorname{Sl}(\underline{n}_1) \arrow[ur,"\beta"'] &
\end{tikzcd}
\]
where \(j\) is the canonical epimorphism and \(\beta\) is a representation of \(\operatorname{G\ell}(\underline{n}_1)/\operatorname{Sl}(\underline{n}_1)\) in \({G}_m\). But the homomorphism
\[
 \det:\operatorname{G\ell}(\underline{n}_i)\longrightarrow{G}_m
\]
has a section, so that the monomorphism
\[
 \operatorname{G\ell}(\underline{n}_i)/\operatorname{S\ell}(\underline{n}_i)\longrightarrow{G}_m
\]
induced by det is an isomorphism. On the other hand, it is known (SGA 3 VIII 1.3) that every representation of \({G}_m\) in itself is of the form
\[
 s\rightsquigarrow s^{\underline{m}},
\]
where \(\underline{m}\) is a section of \(\mathbb Z\), whence the assertion.

Corollary 6.2.6 is a simple reformulation of 6.2.3. Note that 6.2.3 also follows immediately from SGA 3 XXII 6.2.1, but we wish to give an elementary proof here via 6.2.2.

%%%%%%%%%%  scan idx 569  |  folio 556  %%%%%%%%%%
\SGAPrintPage{556}%

\sgaStableTarget{sga6.e10.sec.6.4}\subsection*{6.4. Proof of 6.2.2.}
It suffices to show that, for every \(Y\in\operatorname{Ob}(\Sch/S)\), \(y\in Y\) and \(s\in\operatorname{S\ell}(\underline{n})(Y)\), one can find an ``étale neighborhood'' \(Y'\to Y\) of \(y\) such that the restriction of \(s\) to \(Y'\) is a product of commutators of \(\operatorname{G\ell}(\underline{n})(Y')\). Since the question is local, one may suppose that \(Y\) is affine with ring \(A\), and that \(\underline{n}\) is constant with value \(n\) on \(Y\). Let \(\underline{p}\) be the prime ideal of \(y\) and \((s_{ij})(i,j=1,\ldots,n)\) the matrix of \(s\). For every \(\lambda\in A\) and \(h, k=1,\ldots,n\); \(h\ne k\), we shall denote by \(c_{h,k,\lambda}\) the matrix \((c_{ij})\) \((i,j=1,\ldots,n)\) such that \(c_{hk}=\lambda\) and \(c_{ij}=\delta_{ij}\) for \((i,j)\ne(h,k)\) \((\delta_{ij}=\) the Kronecker symbol). For every matrix \(\underline{a}=(a_{ij})(i,j=1,\ldots,n)\), one obtains \(c_{h,k,\lambda}\underline{a}\) by adding to the \(k\)-th column of \(\underline{a}\) \(\lambda\) times the \(h\)-th column. Similarly, one obtains the matrix \(\underline{a}c_{h,k,\lambda}\) by adding to the \(h\)-th row of \(\underline{a}\) \(\lambda\) times the \(k\)-th row. In particular, one has
\[
 \det(c_{h,k,\lambda})=1,
\]
and one sees at once that
\sgaStableTarget{sga6.e10.eq.6-4-1}\begin{equation}\tag{6.3.1}
 c_{h,k,\lambda}c_{h,k,\mu}=c_{h,k,\lambda+\mu}
\end{equation}
for every \(\lambda,\mu\in A\), so that one obtains a representation
\[
 c_{h,k}:G_a\longrightarrow\operatorname{S\ell}(n)
 \qquad(h,k=1,\ldots,n\ ;\ h\ne k).
\]
Unless expressly stated otherwise, in what follows we shall suppose that the \(\lambda\) occurring in \(c_{h,k,\lambda}\) \((\lambda\in A\ ;\ h,k=1,\ldots,n\ ;\ h\ne k)\) is invertible.
The proof then proceeds in two steps. We shall first show that \(s\) is a product of matrices \(c_{h,k,\lambda}\) in a sufficiently small Zariski neighborhood of \(y\). We shall then show that the matrices \(c_{h,k,\lambda}\) are products of commutators of \(\operatorname{G\ell}(n)\) in a sufficiently small étale neighborhood of \(y\).

%%%%%%%%%%  scan idx 570  |  folio 557  %%%%%%%%%%
\SGAPrintPage{557}%

For the first step, one must transform \(s\) into the identity matrix \(\underline{\underline{1}}=(\delta_{ij})\) in a Zariski neighborhood of \(y\) by multiplying \(s\) on the left and on the right by suitable matrices \(c_{h,k,\lambda}\). Let us first show that one can thereby make \(s_{11}=1\). If \(s_{11}\ne1\) and \(s_{21}\in\underline{p}\), we seek a matrix \(c=c_{2,h,\lambda}\) such that the product \(s'=sc=(s'_{ij})\) has \(s'_{21}\notin\underline{p}\). Since \(\det(s)=1\notin\underline{p}\), there is at least one index \(i\) such that \(s_{i1}\notin\underline{p}\). Then one has \(s_{21}+s_{i1}\notin\underline{p}\), and \(c_{2,i,1}\) answers the question. One may therefore suppose \(s_{21}\notin\underline{p}\) and, after localizing, suppose that \(s_{21}\) is invertible. If \(s_{11}-1\in\underline{p}\), then in \(s'=(s'_{ij})=sc_{1,2,1}\) one has
\(s'_{11}-1=(s_{11}+s_{21})-1\notin\underline{p}\). After localizing, one may then suppose that \(s_{11}-1\) is invertible. Then \(\lambda=(1-s_{11})/s_{21}\) is invertible and one finds \(s'_{11}=1\) in \(s'=(s'_{ij})=sc_{1,2,\lambda}\).

Supposing \(s_{11}=1\), one can then put \(s\) in the form
\[
 s_{11}=1,\qquad s_{1i}=s_{i1}=0\quad(i=2,\ldots,n)
\]
in a Zariski neighborhood, of \(y\) by multiplying \(s\) on the left (resp. on the right) by matrices \(c_{1,i,\lambda}\) (resp. \(c_{i,1,\lambda}\)). Next one makes \(s_{22}=1\), etc. By induction, one can then transform \(s\) in a Zariski neighborhood \(U\) of \(y\) into a matrix \(s'=(s'_{ij})\) with \(s'_{ij}=\delta_{ij}\) for \(i<n\) or \(j<n\). But by the very construction one then has \(s'=c_1sc_2\) on \(\mathcal U\), where \(c_1\) and \(c_2\) are products of matrices \(c_{h,k,\lambda}\). Thus one has
\[
 \det(s')=\det(s)=1,
\]
so that \(s'_{nn}=1\), \(s'=\underline{\underline{1}}\), and the result follows.

We therefore suppose that \(s\) is already, on \(Y\), a product of matrices \(c_{h,k,\lambda}\).
It remains to see that these are products of commutators in an étale neighborhood of \(y\). First, one sees that any two matrices

%%%%%%%%%%  scan idx 571  |  folio 558  %%%%%%%%%%
\SGAPrintPage{558}%

\(c_{h,k,\lambda}\) and \(c_{h,k,\mu}\) are conjugate; in other words, there is an invertible matrix \(a=(a_{ij})\) \((i,j=1,\ldots,n)\) such that
\[
 ac_{h,k,\lambda}=c_{h,k,\mu}a.
\]
Indeed, the matrix \((a_{ij})\) with \(a_{hh}=1/\mu\), \(a_{kk}=1/\lambda\), and otherwise \(a_{ij}=\delta_{ij}\), answers the question (recall that \(\lambda\) and \(\mu\) are invertible by hypothesis and \(h\ne k\)). Thus, for any homomorphism \(\beta\) from \(\operatorname{G\ell}(n)(U)\) to a sheaf of abelian groups \(H\), one has
\[
 \beta(c_{h,k,\lambda})=\beta(c_{h,k,\mu}).
\]
If \(\lambda+\mu\) is invertible, this and (6.3.1) imply that \(\beta(c_{h,k,\lambda})=0\) for every \(h,\beta\), so that \(c_{h,k,\lambda}\) is a product of commutators. If the residue field \(k\) of \(y\) is not the prime field \(\underline{F}_2=\{0,1\}\), then for every invertible \(\lambda\in A\) one can find a Zariski neighborhood \(U=\Spec(B)\) of \(y\) and an invertible element \(\mu\in B\) such that \(\lambda+\mu\in B\) is invertible. It follows that \(c_{h,k,\lambda}\) is a product of commutators on \(U\).

If \(k=\underline{F}_2\), then \(V=\Spec(A[T]/(T^2+T+1))\) is an étale neighborhood of \(y\) whose residue field differs from \(\underline{F}_2\). The result follows.

\sgaStableTarget{sga6.e10.sec.6.5}\subsection*{6.5. Remarks}
\sgaStableTarget{sga6.e10.para.6-5-1}\textbf{6.5.1.} When none of the residue fields of \(S\) is isomorphic to \(\underline{F}_2\), the preceding proof shows that 6.2.2. remains valid for the ringed topos of \emph{Zariski} sheaves on \(S\). On the other hand, already for \(S=\Spec(\underline{F}_2)\), \(\operatorname{S\ell}(2,\underline{F}_2)=\operatorname{G\ell}(2,\underline{F}_2)\) differs from the commutator group of \(\operatorname{G\ell}(2,\underline{F}_2)\). The order of the latter is \(3\), whereas the order of \(\operatorname{S\ell}(2,\underline{F}_2)\) is \(6\) (cf. [1]).

\sgaStableTarget{sga6.e10.para.6-5-2}\textbf{6.5.2.} On an arbitrary locally ringed topos \(X\), it is not true in general that every representation \(G_m\to G_m\) is defined by a
